\section{Uniform indices in closed affine groups}

Fix a prime $p$ and write $R=\mathbb Z_p$. The affine multiplication
is $(b,u)(c,v)=(b+uc,uv)$ on $R\rtimes R^\times$. Every subgroup
and index below is an actual subgroup or coset index; an assertion
that an index is finite refers to the corresponding coset set.

\begin{lemma}[A uniform scalar index]
\label{ft:uniform-scalar-index}
\uses{ft:padic-closed-additive-subgroups}
\lean{PadicInt.mulLeft_range_finiteIndex, PadicInt.mulLeft_range_index_pos, PadicInt.closedAddSubgroup_map_mulLeft_relIndex_eq, PadicInt.closedAddSubgroup_map_mulLeft_relIndex_le}
\leanok
Let $a\in R$ be nonzero, and put $C(a)=[R:aR]$. This is a positive
finite integer. For every closed additive subgroup $T$ of $R$,
\[
 [T:aT]\leq C(a),
\]
with equality if $T\ne0$ and index one if $T=0$.
\end{lemma}
\begin{proof}
The ideal $aR$ is nonzero, hence open in the compact discrete valuation
ring $R$. Its quotient is discrete and compact, so finite; its zero
coset makes its cardinal positive. Closed additive subgroups of $R$
are ideals by density of the integers and continuity of scalar
multiplication. If $T\ne0$, their classification gives $T=p^rR$ for
some $r\geq0$. Multiplication by $p^r$ is an additive isomorphism
$R\longrightarrow T$: it is surjective by this description and
injective because $R$ is a domain. This isomorphism sends $aR$
exactly onto $aT$, and therefore induces a bijection
$R/aR\longrightarrow T/aT$. This proves equality of the indices.
If $T=0$, its quotient has one element; positivity of $C(a)$ gives
the claimed inequality also in this case.
\end{proof}

\begin{lemma}[The zero stabilizer modulo the translation kernel]
\label{ft:affine-zero-stabilizer}
\uses{ft:padic-closed-additive-subgroups,ft:affine-commutator}
\lean{FalseTateExtension.affinePureUnits, FalseTateExtension.mem_affinePureUnits, FalseTateExtension.isClosed_affinePureUnits, FalseTateExtension.affineProjection_injective_on_pureUnits, FalseTateExtension.affinePureUnits_commute, FalseTateExtension.affineZeroStabilizer, FalseTateExtension.mem_affineZeroStabilizer, FalseTateExtension.affineZeroStabilizer_eq_sup, FalseTateExtension.affineZeroStabilizer_inv_mul_mem_iff, FalseTateExtension.affineZeroStabilizerQuotientEmbedding, FalseTateExtension.affineZeroStabilizer_finiteIndex, FalseTateExtension.affinePureUnits_image_index, FalseTateExtension.affinePureUnits_image_finiteIndex, FalseTateExtension.exists_affinePureUnit_ne_one}
\leanok
Let $G$ be a closed subgroup of $R\rtimes R^\times$, let
$\chi:G\longrightarrow H\subseteq R^\times$ be its actual
surjective unit projection, and identify its full translation kernel
with
\[
 M=\{b\in R:(b,1)\in G\}.
\]
Put $S=G\cap(\{0\}\rtimes R^\times)$. The groups $M$ and $S$ are
closed, $S$ is abelian, and $M\cap S=1$. If $M\ne0$, then
\[
 D=[G:MS]=[H:\chi(S)]
\]
is positive and finite. If in addition $H$ is infinite, $S$ contains
an element $s_0$ with $\chi(s_0)\ne1$.
\end{lemma}
\begin{proof}
The coordinate inclusions and projections are continuous; the defining
closed conditions give closedness of $M$ and $S$. The multiplication
on $S$ is multiplication of its unit coordinate, so $S$ is abelian
and $\chi$ is injective on $S$. Its intersection with the translation
kernel is consequently trivial.

Assume $M\ne0$. The closed-additive-subgroup classification gives
$M=p^rR$ for some $r\geq0$. In particular $R/M$ is finite: the ideal
$M$ is open and $R$ is compact. Since $M$ is an ideal, each affine
map $x\mapsto b+ux$ for $(b,u)\in G$ preserves congruence modulo
$M$. The inverse affine map does the same. The affine multiplication
formula therefore gives an action of $G$ on the finite set $R/M$.
The stabilizer of zero consists exactly of elements with $b\in M$.
For such an element, $(b,1)\in G$ and
$(b,1)^{-1}(b,u)=(0,u)\in S$; conversely every element of $MS$
has translation coordinate in $M$. Thus the stabilizer is $MS$.
Its cosets are in bijection with the orbit of zero, proving
$1\leq D\leq |R/M|<\infty$.

The full kernel of $\chi$ is $M$, so the inverse image of $\chi(S)$
is exactly $MS$. Surjectivity of $\chi$ induces a bijection of
coset sets $G/MS\longrightarrow H/\chi(S)$, proving their index
equality. If $H$ is infinite, its finite-index subgroup $\chi(S)$
is infinite: a finite subgroup with finitely many cosets would make
$H$ finite. It thus contains an element different from one, which
lifts to the required $s_0\in S$.
\end{proof}

\begin{theorem}[Uniform comparison of the two projection indices]
\label{ft:closed-affine-index-comparison}
\uses{ft:uniform-scalar-index,ft:affine-zero-stabilizer,ft:affine-commutator}
\lean{FalseTateExtension.affineTranslationIntersection, FalseTateExtension.affineTranslationIntersection_toAddSubgroup, FalseTateExtension.affineTranslationIntersection_relIndex, FalseTateExtension.isClosed_affinePureUnits_sup, FalseTateExtension.affinePureUnit_commutator_mem, FalseTateExtension.affineTranslationIntersection_mono, FalseTateExtension.affinePureUnit_scalar_image_le, FalseTateExtension.affinePureUnit_intersection_relIndex_bound, FalseTateExtension.affinePureUnit_kernel_index_bound, FalseTateExtension.affineIndex_finiteIndex_of_projectionIndices, FalseTateExtension.affineIndex_comparison}
\leanok
Keep the actual closed affine group $G$, full translation kernel $M$,
unit image $H$, and pure-unit subgroup $S$ of the preceding lemma.
Assume $M\ne0$ and $H$ is infinite. There are positive integers
$C,D$, depending only on $G$, with the following property. For every
closed normal subgroup $I$ of $G$ such that
\[
 A=[G:SI],\qquad B=[H:\chi(I)]
\]
are finite, the quotient $G/I$ is finite, $I$ is open, and
\[
 AB\leq D[G:I],\qquad [G:I]\leq CAB.
\]
Equivalently $AB/D\leq[G:I]\leq CAB$ as real inequalities.
No openness of $I$ is assumed.
\end{theorem}
\begin{proof}
Take $D=[G:MS]$ and choose $s_0\in S$ with
$\chi(s_0)\ne1$ by the preceding lemma. Put
$a=\chi(s_0)-1\ne0$ and $C=C(a)$. These choices are independent
of $I$, and both constants are positive and finite.

For a given $I$ put $M_I=M\cap I$ and $T=M\cap SI$,
identifying translations with their coordinates in $R$. Normality
of $I$ makes $SI$ a subgroup. The group $G$ is compact, and both
$S$ and $I$ are closed. The image of the compact set $S\times I$
under multiplication is $SI$, so $SI$ and then $T$ are closed.
For $t\in T$, write $t=si$ with $s\in S$ and $i\in I$.
Modulo $I$, the element $t$ is represented by $s$. Since $S$ is
abelian, $s_0$ commutes with this class, whence $[s_0,t]\in I$.
The actual affine commutator formula identifies this commutator
with the translation $at$. Thus
\[
 aT\subseteq M_I\subseteq T.
\]
The scalar-index lemma makes $T/aT$ finite of size at most $C$.
The quotient map from its coset set onto $T/M_I$ gives
$[T:M_I]\leq C$ and finiteness of $T/M_I$.

Inclusion of $M$ into $G$ induces an injection of coset sets
$M/T\longrightarrow G/SI$: two elements of $M$ have the same
image exactly when their quotient is in $M\cap SI=T$.
Consequently $M/T$ is finite and $[M:T]\leq A$. The index
formula for $M_I\subseteq T\subseteq M$ gives
\[
 [M:M_I]=[M:T][T:M_I]\leq CA,
\]
so this index is finite as well. To obtain the reverse comparison,
we first establish finiteness of the full quotient. The surjection
$G/I\longrightarrow H/\chi(I)$ has kernel
$MI/I$, and the latter is isomorphic to $M/M_I$. These assertions
follow by choosing an element of $I$ with the same unit coordinate
as any element over $\chi(I)$. Thus both kernel and target are
finite, which proves finiteness of $G/I$ and the exact formula
\[
 [G:I]=B[M:M_I].
\]
Write $Q=G/I$, and let $N,V$ be the respective images of $M,S$ in
this finite group. Coset correspondence gives
\[
 [Q:V]=A,\quad |N|=[M:M_I],\quad
 [Q:NV]=[G:MSI]\leq[G:MS]=D.
\]
The subgroup $N$ is normal, so $NV$ is the image of $N\times V$
under multiplication. Consequently $|NV|\leq|N||V|$. The finite
group cardinality formulas now give
\[
 A|V|=|Q|=[Q:NV]|NV|\leq D|N||V|.
\]
The group $V$ contains the identity, so its finite cardinal is
positive. Cancellation proves $A\leq D[M:M_I]$. Multiplication
of this inequality and the already proved $[M:M_I]\leq CA$ by
$B$ gives both stated bounds. Finally $I$ is closed with finitely many cosets
in $G$; its complement is the finite union of the other closed
cosets. Hence $I$ is open.
\end{proof}

\begin{lemma}[Zero translation kernel]
\label{ft:affine-zero-kernel-isomorphism}
\uses{ft:affine-commutator}
\lean{FalseTateExtension.affineProjection_ker_eq_bot_of_translationKernel_eq_bot, FalseTateExtension.affineProjection_injective_of_translationKernel_eq_bot, FalseTateExtension.affineProjectionEquivOfTranslationKernelEqBot, FalseTateExtension.affineProjectionEquivOfTranslationKernelEqBot_apply, FalseTateExtension.affineProjection_index_of_translationKernel_eq_bot}
\leanok
For a closed affine group $G$ with zero full translation kernel,
the actual unit projection $\chi:G\longrightarrow H$ is an
isomorphism of topological groups. For every subgroup $I$ of $G$
it induces a bijection $G/I\longrightarrow H/\chi(I)$, preserving
the actual index.
\end{lemma}
\begin{proof}
The unit projection is surjective onto its actual image by definition.
Its kernel is the zero translation subgroup, so it is injective.
It is continuous. The affine group is compact and Hausdorff, so
its closed subgroup $G$ is compact, while $H$ is Hausdorff as a
subspace of $R^\times$. A continuous bijection from a compact
space to a Hausdorff space is a homeomorphism. The projection and
its inverse are group homomorphisms, proving the assertion.
The same isomorphism maps $I$ onto $\chi(I)$ and hence identifies
their coset sets.
\end{proof}
